\documentclass{article}
\usepackage[utf8]{inputenc}
\usepackage[T1]{fontenc}
\usepackage[a4paper]{geometry}		% Reduire les marges
\geometry{hmargin=3cm,vmargin=3.5cm}
\usepackage[english]{babel} 
\usepackage{amsfonts}
\usepackage{amsthm}
\usepackage{amsmath}
\usepackage{amssymb}
\usepackage{hyperref}
\usepackage{array}
\usepackage[svgnames]{xcolor}
\usepackage{xpatch}
\usepackage{tikz}
\usepackage{mathrsfs}
%\usepackage{graphicx}
%\usepackage{caption}
\usepackage[svgnames]{xcolor}
\usepackage{eurosym}
\usetikzlibrary{positioning,calc}
\renewcommand{\thesection}{\Roman{section}}
\xpatchcmd{\proof}{\itshape}{\normalfont\proofnameformat}{}{}
\newcommand{\proofnameformat}{\bfseries}
\theoremstyle{definition}
\newtheorem{defi}{Definition}[section]
\newtheorem{theo}{Theorem}[section]
\newtheorem{prop}[theo]{Proposition}
\newtheorem{lem}[theo]{Lemma}
\newtheorem{coro}[theo]{Corollary}
\newtheorem{exe}{Example}[section]
\newtheorem{rem}{Remark}[section]
\newtheorem{nota}{Notation}[section]
\renewcommand\qedsymbol{$\blacksquare$}
\numberwithin{equation}{section}

\DeclareMathOperator{\R}{\mathbb{R}}
\DeclareMathOperator{\Z}{\mathbb{Z}}
\DeclareMathOperator{\Q}{\mathbb{Q}}
\DeclareMathOperator{\N}{\mathbb{N}}
\DeclareMathOperator{\E}{\mathbb{E}}
\DeclareMathOperator{\C}{\mathcal{C}}

\title{Exo 6, corrigé}       
\author{Guillaume Woessner}
\date{}                       % sans date : celle du jour de compilation

%\pagestyle{headings}          % Pour mettre des entetes avec les titres
                              % des sections en haut de page
\sloppy                       % Ne pas faire deborder les lignes dans la marge

\begin{document}

\maketitle                    % Faire un titre utilisant les donnees
                              % passees : \title, \author et \date

\begin{abstract}
\begin{center}
Les modes de convergence
\end{center}
\end{abstract}

%\tableofcontents              % Table des matieres


Dans tout ce document, l'espace de probabilité filtré considéré est $(\Omega,\mathcal{F},(\mathcal{F}_n)_n,\mathbb{P})$, et on supposera que $X$ est une variable aléatoire, et que $(X_n)_n$ est une suite de variables aléatoires adaptée à la filtration.


\section{Recap}

\paragraph{Exercice 1}
Montrez avec des contre-exemples que les implications suivantes sont fausses :
\begin{align*}
    X_n \mapsto_n X \text{ en } \mathcal{L} ~~ &\nRightarrow~~  X_n \mapsto_n X \text{ en } \mathbb{P} & (A)\\
    X_n \mapsto_n X \text{ en } \mathbb{P} ~~ &\nRightarrow~~  X_n \mapsto_n X \textit{ p.s.}& (B)\\
    X_n \mapsto_n X \text{ en } \mathbb{P} ~~ &\nRightarrow~~  X_n \mapsto_n X \text{ dans } L^p ~~~ (p\geq 1)& (C)\\
    X_n \mapsto_n X \textit{ p.s.} ~~ &\nRightarrow~~  X_n \mapsto_n X \text{ dans } L^p ~~~ (p\geq 1)& (D)\\
    X_n \mapsto_n X \text{ dans } L^p ~~~ (p\geq 1) ~~ &\nRightarrow~~  X_n \mapsto_n X \textit{ p.s.}& (E)
\end{align*}

\begin{proof}~
\begin{itemize}
\item[(A)] On peut poser $Y$ suivant une loi de Rademacher, et $X_{2n}=Y$ et $X_{2n+1}=-Y$.
\item[(B)] On se place en particulier sur l'espace de probabilité $([0,1],\mathcal{B},m)$ avec $m$ la mesure de Lebesgue, et on construit la suite de variables aléatoires $X_0:=a_0\mathbf{1}_{[0,1[}$, $X_1:=a_1\mathbf{1}_{[0,1/2[}$, $X_2:=a_2\mathbf{1}_{[1/2,1[}$, $X_3:=a_3\mathbf{1}_{[0,1/3[}$, $X_4:=a_4\mathbf{1}_{[1/3,2/3[}$, $X_5:=a_5\mathbf{1}_{[2/3,1[}$, $X_6:=a_6\mathbf{1}_{[0,1/4[}$, etc...\\
Poser $a_n=1$ fourni un contre-exemple avec $X=0$.
\item[(C)] En posant $a_n = l(n)^{-1/p}$ où $l(n)$ est la longueur de l'intervalle non nul au rang $n$, le même contre-exemple fonctionne.
\item[(D)] Si les $X_n$ ou bien $X$ ne sont pas dans $L^p$, il ne peut y avoir convergence dans $L^p$.
\item[(E)] Le même contre-exemple que pour (B) fonctionne.
\end{itemize}
\end{proof}


\paragraph{Exercice 2}
Supposons que $X_n \mapsto_n X \text{ en } \mathcal{L}$, et que $X$ soit une constante $c$ presque sûrement.\\
Montrez que dans ce cas $X_n \mapsto_n X \text{ en } \mathbb{P}$.

\begin{proof}
Pour des raisons de simplicité, on suppose que $X=c=0$, mais la preuve s'adapte facilement au cas général.\\
En écrivant $F_n$ la fonction de répartition de $X_n$ et $F$ celle de $X$ (donc $F(x)=\mathbf{1}_{x\geq 0}$) on a, pour $\varepsilon>0$,
\begin{align*}
\mathbb{P}(|X_n|>\varepsilon) &= \mathbb{P}(X_n<-\varepsilon) + \mathbb{P}(X_n>\varepsilon) = F_n(-\varepsilon) + 1 - F_n(\varepsilon)\\
&\mapsto_n F(-\varepsilon) + 1 - F(\varepsilon) = 0 + 1 - 1 = 0.
\end{align*}
\end{proof}


\paragraph{Exercice 3}
Supposons que 
$$ \sum_n \mathbb{P}(|X_n-X|>\varepsilon) < \infty. $$
Montrez que $X_n \mapsto_n X \text{ en } \mathbb{P}$, et (plus difficile), montrez que $X_n \mapsto_n X \textit{ p.s.}$.\\
Déduisez-en que si $X_n\mapsto X$ en $\mathbb{P}$, il existe une sous-suite qui converge presque sûrement.

\begin{proof}
Pour la convergence en probabilité, il suffit de remarquer que la condition implique que $\mathbb{P}(|X_n-X|>\varepsilon)\mapsto_n 0$.\\
Pour la convergence presque sûre il suffit d'appliquer le lemme de Borel-Cantelli aux évènements $\{|X_n-X|>\varepsilon\}$.\\
Pour le corollaire, par hypothèse, on sait que pour tout $k>0$ il existe un rang $N_k$ à partir duquel pour tout $n\geq N_k$ on a $\mathbb{P}(|X_n-X|>\varepsilon)<2^{-k}$. On peut donc construire une suite croissante $(n_k)_k$ telle que $\mathbb{P}(|X_{n_k}-X|>\varepsilon)<2^{-k}$, et ainsi pouvoir appliquer ce qui précède à la suite $(X_{n_k})_k$.
\end{proof}


\section{Cours}


\paragraph{Exercice 5}
Supposons que 
$$\forall \varepsilon>0,~\forall\delta>0,~\exists n\in\mathbb{N} \textit{ s.t. }~~\mathbb{P}(\sup_{m\geq n} |X_n-X_m|>\varepsilon)<\delta.$$
Montrez que $X_n \mapsto_n X \text{ en } \mathbb{P}$, et (plus difficile), montrez que $X_n \mapsto_n X \textit{ p.s.}$.

\begin{proof}
Pour la convergence en probabilité, il suffit de remarquer que $\mathbb{P}(\sup_{m\geq n} |X_n-X_m|>\varepsilon)\geq \sup_{m\geq n} \mathbb{P}(|X_n-X_m|>\varepsilon)$.\\
Pour la convergence presque sûre, en reformulant l'hypothèse on sait que, pour tout $\varepsilon>0$,
$$\lim_n \mathbb{P}(\sup_{m\geq n}|X_n-X_m|>\varepsilon)=0.$$
Ainsi, d'après le lemme de Fatou,
\begin{align*}
0&= \lim_n \mathbb{P}(\sup_{m\geq n}|X_n-X_m|>\varepsilon) = \liminf_n \mathbb{P}(\sup_{m\geq n}|X_n-X_m|>\varepsilon) \\
&\geq \mathbb{P}(\liminf_n \{ \sup_{m\geq n}|X_n-X_m|>\varepsilon\}) = \mathbb{P}(\cup_k \cap_{n\geq k} \{ \sup_{m\geq n}|X_n-X_m|>\varepsilon\})\\
&= \mathbb{P}(\cup_k \cap_{n\geq k}  \cup_{m\geq n} \{|X_n-X_m|>\varepsilon\}) \geq \mathbb{P}( \cap_n  \cup_{m\geq n} \{|X_n-X_m|>\varepsilon\})\\
&= \mathbb{P}(\limsup \{|X_n-X_m|>\varepsilon\})
\end{align*}
QED
\end{proof}




\paragraph{Exercice 6}[Doob]~\\
Soient des constantes $1<p<\infty$ et $a>0$, soit $(X_n)_n$ une sous-martingale positive et le temps d'arrêt $\tau_a:=\inf\{n\geq 0~:~X_n>a\}$. On note $S_n:=\sup_{k\leq n}X_k$.\\
Montrez la version plus fine de l'\textbf{inégalité maximale de Doob} suivante,
$$\mathbb{P}(S_n>a)\leq \dfrac{1}{a}\E[X_n\mathbf{1}_{\tau_a\leq n}].$$
\textit{Indice :} Remarquez que pour tout $ k\leq n$, on a : $a\mathbf{1}_{\tau_a =k} \leq \E[X_n~|~\mathcal{F}_k]\mathbf{1}_{\tau_a =k}$.
~\\
Utilisez ce résultat pour montrer l'\textbf{inégalité $L^p$ de Doob} suivante,
$$\E[S_n^{~p}]\leq \left(\dfrac{p}{p-1}\right)^p \E[X_n^{~p}].$$

\begin{proof}
On a, en suivant l'indication
$$a\mathbf{1}_{\tau_a =k} \leq X_k\mathbf{1}_{\tau_a =k} \leq  \E[X_n~|~\mathcal{F}_k]\mathbf{1}_{\tau_a =k}.$$
Ainsi on peut calculer, en utilisant le point précédent,
\begin{align*}
\mathbb{P}(S_n>a)&=\mathbb{P}(\tau_a\leq n) = \sum_{k=0}^n\mathbb{P}(\tau_a=k)\\
&=\sum_{k=0}^n\mathbf{1}_{\tau_a=k} \leq \dfrac{1}{a}\sum_{k=0}^n\E\big[\E[X_n~|~\mathcal{F}_k]\mathbf{1}_{\tau_a=k}\big]\\
&=\dfrac{1}{a}\sum_{k=0}^n\E\big[\E[X_n\mathbf{1}_{\tau_a=k}~|~\mathcal{F}_k]\big]=\dfrac{1}{a}\sum_{k=0}^n\E[X_n\mathbf{1}_{\tau_a=k}]\\
&=\dfrac{1}{a}\E[X_n\mathbf{1}_{\tau_a\leq n}]
\end{align*}
On effectue alors le calcul suivant, en notant $q$ le conjugué de $p$, c'est-à-dire le nombre tel que $\frac{1}{q}+\frac{1}{p}=1$,
\begin{align*}
\E[S_n^p]&=\int_0^\infty \mathbb{P}(S_n>a^{1/p})da \leq \int_0^\infty \dfrac{1}{a^{1/p}}\E[X_n\mathbf{1}_{\tau_{a^{1/p}}\leq n}]da\\
&=\E\left[X_n\int_0^\infty \dfrac{1}{a^{1/p}}\mathbf{1}_{\tau_{a^{1/p}}\leq n}da\right] = \E\left[X_n\int_0^\infty \dfrac{1}{a^{1/p}}\mathbf{1}_{S_n^{~p}>a}da\right]\\
&= \E\left[X_n\int_0^{S_n^{~p}} \dfrac{1}{a^{1/p}}da\right] = \E\left[X_n \dfrac{(S_n^{~p})^{1-\frac{1}{p}}}{1-\frac{1}{p}} \right]\\
&= \dfrac{p}{p-1}\E\left[X_n S_n^{~p-1} \right].
\end{align*}
Enfin, appliquer l'inégalité de Hölder avec $p$ et $q$ permet de conclure.
\end{proof}



\paragraph{Exercice 7} Vérifiez que la preuve du théorème de convergence des martingales dans $L^2$ reste valide si on suppose seulement que $(X_n)_n$ est une sous-martingale.\\
En déduire que le théorème est vraie pour une sur-martingale également.\\
\textbf{Attention !} la convergence n'a plus lieu que presque sûrement, plus dans $L^2$.

\begin{proof}
La preuve s'effectue de la même manière, le lemme préliminaire devient seulement :\\
Pour tout $m\geq n$, on a 
$$\E[(X_m-X_n)^2]\leq \E[X_m^2]-\E[X_n^2].$$
Ensuite, il suffit de noter que l'opposée d'une sur-martingale est une sous-martingale.\\
On peut remarquer que la convergence n'a plus lieu presque sûrement en considérant la suite de variables aléatoires de $[0,1]$ à $\R$ définies par $X_n(x) := 2^{-n/2}\mathbf{1}_{[0,2^{-n}]}(x)$.
\end{proof}


















\end{document}












